\documentclass[12pt]{article}

\usepackage{comment}
\usepackage{amsmath}
\usepackage{amssymb}  % for \nmid

\newcommand{\BILL}{{\rm BILL}}
\newcommand{\D}{{\mathbb D}}
\newcommand{\G}{{\mathbb G}}
\newcommand{\Z}{{\mathbb Z}}
\newcommand{\N}{{\mathbb N}}
\newcommand{\Q}{{\mathbb Q}}
\newcommand{\R}{{\mathbb R}}
\newcommand{\succc}{{\rm succ}}
\newcommand{\pred}{{\rm pred}}

\newcommand{\lc}{\left\lceil}
\newcommand{\rc}{\right\rceil}
\newcommand{\Ceil}[1]{\left\lceil {#1}\right\rceil}
\newcommand{\ceil}[1]{\left\lceil {#1}\right\rceil}
\newcommand{\floor}[1]{\left\lfloor{#1}\right\rfloor}

\begin{document}


\centerline{Take Home Part of the Final DUE May 12 at 10:30AM- NO DEAD CAT}

\newif{\ifshowsoln}
\showsolntrue % comment out to NOT show solution inside \ifshowsoln and \fi blocks.

\begin{enumerate}
\item
(0 points)
What is your name.
\textbf{Gurkeerat Singh}

\vfill 

\centerline{\bf GO TO THE NEXT PAGE}

\newpage

\item 
(20 points on the final) 

{\bf For PART $a$ of this problem we want the answer in terms of
binomial coefficients, powers, $\times$, $+$, $-$, division, floor, ceiling. 
For example, an answer like $\floor{\frac{{2^{10}\times 5-10}}{\binom{9}{2}\times 88}}$
(this is NOT the answer) 
would be in the right form, whereas an answer like 286,000
would NOT be in the right form.} 


\bigskip
\bigskip

And now for the problem!

\bigskip

\bigskip

Let $\BILL$ be the following operation on
sets $\{x_1,\ldots,x_{10}\}$ of ten distinct natural numbers: 

$\BILL(x_1,\ldots,x_{10})=x_1^2 + \cdots x_{10}^2$. 

\begin{enumerate}
\item 
(Show all work.) 
Fill in the $n$ in the following sentence and prove the result:

{\it Let $A\subseteq \{1,\ldots,1000\}$  of size 50. 
Then there exists $n$ sets

$$A_1,\ldots,A_n$$

\noindent
such that:
\begin{itemize}
\item 
For all $i$, $|A_i|=10$. 
\item 
$\BILL(A_1)=\BILL(A_2)=\cdots= \BILL(A_n)$. 
\end{itemize} 
}
\begin{itemize}
    \item \textbf{So, first we know that there are $50 \choose 10$ ways to choose a set $|A_i|$ of 10 elements from a set $A$ of 50 elements. Now, we have to find the possible end sums for calling the function BILL on our set $A_i$. The lower bound is if we assume that the lowest possible numbers, $1,2,...,10$, are chosen. The sum would be $1^2 + 2^2 + \cdots + 10^2 = 385$. Further, the upper bound for our sum is if we assume the highest possible numbers, $991,992,...,1000$, are chosen. The sum would be $991^2 + 992^2 + \cdots + 1000^2$, equal to 9910285. So, we have a range of possible sums from 385 to 9910285, meaning that there are $9910285-385+1 = 9909901$ possible sums. So, by using the pigenhole principle, we can use $\left\lceil \frac{\binom{50}{10}}{9909901} \right\rceil$ to determine the number of sets n which meet the requirement.}
\end{itemize}

\item
NOW give me the answer as an actual number. 
\begin{itemize}
    \item \textbf{The answer is $\left\lceil \frac{\binom{50}{10}}{9909901}\right\rceil = 1037$.}
\end{itemize}

\end{enumerate}

\end{enumerate}

\end{document}
